\section{Experimental Setup}
\label{sec:setup}

\subsection{Datasets}

We select two datasets to create a deliberate contrast between domain-matched
and domain-mismatched retrieval, testing whether our findings depend on the
encoder-corpus alignment quality.

\textbf{SQuAD}~\citep{rajpurkar2016squad} contains Wikipedia-sourced evidence
paragraphs paired with reading comprehension questions. The MiniLM encoder
was pre-trained on general web text, making this a \emph{domain-matched}
setting where we expect retrieval to be fundamentally sound. This lets us
isolate the effect of refinement decisions on an otherwise strong pipeline.

\textbf{PubMedQA} \citep{jin2019pubmedqa} contains biomedical research
abstracts with yes/no/maybe questions about clinical findings. Applying the
same general-domain encoder to specialized scientific language creates a
\emph{domain-mismatched} setting that tests whether coherence-based
insights generalize when upstream retrieval is systemically weak.

This two-dataset design is intentional: any method that helps only on SQuAD
could be an artifact of strong retrieval; any insight that holds on both
datasets is more likely to reflect a general principle about context
coherence.

\subsection{Models}

Generation is performed with Mistral-7B~\citep{jiang2023mistral} and
Llama-3-8B~\citep{llama3}. Both are instruction-tuned and served through the
same local inference stack (Ollama v0.1.29) to ensure that dataset and
retrieval effects are not confounded by deployment differences. Using two
architecturally distinct models tests whether our findings are
model-specific or reflect a more general property of how language models
interact with retrieved context.

\subsection{Phased evaluation design}

The study is organized into six experimental phases, each isolating a
specific variable while reusing the shared pipeline. This structure prevents
confounding and enables direct comparison across ablations.

\begin{table}[!htb]
\centering
\caption{Experimental phases and their role in the study. The phased design
lets us test each variable in isolation while reusing the same pipeline
across all conditions.}
\label{tab:phases}
\small
\begin{tabular}{clcl}
\toprule
Phase & Focus & Queries & Primary variable \\
\midrule
1 & Chunk size ablation    & 720  & chunk $\in \{256, 512, 1024\}$ \\
2 & Multi-model validation & 1440 & Mistral vs.\ Llama-3 \\
3 & Re-ranking study       & 240  & with/without cross-encoder \\
4 & Zero-shot baseline     & 350  & retrieval vs.\ no retrieval \\
5 & Threshold sweep        & 3000 & $\tau_{\text{sim}}$, $\tau_{\text{ce}}$ \\
6 & HCPC head-to-head      & 180  & baseline vs.\ v1 vs.\ v2 \\
\bottomrule
\end{tabular}
\end{table}

\input{figures/study_design}
Figure~\ref{fig:study_design} visualizes the phase structure.
Phase~1 evaluates the effect of chunk size, top-$k$, and prompt strategy
under a fixed generator. Phase~2 replicates the chunking study across both
model families to test generalization. Phase~3 compares cross-encoder
re-ranking against HCPC variants to separate scoring effects from refinement
effects. Phase~4 measures whether retrieval actually helps relative to
zero-shot prompting---the most basic sanity check for any RAG system.
Phase~5 explores threshold sensitivity across 168 configurations to verify
that HCPC-v2 does not depend on brittle parameter choices. Phase~6 is the
most controlled comparison and directly contrasts baseline, HCPC-v1, and
HCPC-v2 under identical conditions.

\subsection{Query sampling, power, and reproducibility}

All experiments use queries sampled randomly from the validation split at the
start of each phase and held fixed across all conditions within that phase.
This ensures that condition comparisons are not confounded by query
difficulty variation.

The study uses a deliberate two-tier sample design. Phases~1--2 use
$n = 100$ per condition (totaling 2,160 queries) for discovery and variance
analysis, providing the statistical power behind the ANOVA results and effect
size estimates. Phase~6, the HCPC head-to-head comparison, uses $n = 30$ per
condition. This is a \emph{confirmatory} phase, not a discovery phase: it
tests whether the trends established across 6,000+ earlier queries replicate
under maximally controlled conditions (same queries, same model, same
corpus, only the retrieval strategy varies). At $n = 30$, a two-sample
$t$-test detects effects of Cohen's $d \geq 0.75$ at $\alpha = 0.05$ with
80\% power. The effects we report are substantially larger: the
baseline-to-HCPC-v1 faithfulness drop on SQuAD corresponds to $d > 1.0$,
well above this detection threshold. We report 95\% bootstrap confidence
intervals for key faithfulness estimates in Section~\ref{sec:results} to
make the precision of Phase~6 estimates transparent.

We acknowledge that Phase~6's per-condition samples are small in absolute
terms. The hallucination rates (reported as percentages of $n = 30$) have a
resolution of 3.3~pp, meaning the 0\% vs.\ 10\% difference on SQuAD
reflects 0 vs.\ 3 hallucinated responses. We do not claim precision at the
individual-query level from Phase~6; its role is to confirm direction and
approximate magnitude of effects established in the larger phases.

Critically, every finding from Phase~6 is consistent with the larger
phases. The chunk-size dominance (90.1\% of between-factor variance) is
estimated from 720 queries. The model-specific variance decomposition uses
1,440 queries. The threshold robustness result spans 3,000 configurations.
Phase~6 confirms these trends under the most controlled comparison rather
than standing alone as the sole source of evidence.

\subsection{Implementation environment}

All experiments were run on Apple Silicon hardware with MPS acceleration for
embedding and NLI inference. ChromaDB v0.4 served as the persistent vector
store. Because all conditions share the same infrastructure, runtime
differences are attributable to retrieval and refinement behavior rather
than to heterogeneous serving backends. The full codebase is implemented
in Python 3.11 using \texttt{langchain}, \texttt{sentence-transformers},
\texttt{chromadb}, and \texttt{torch} with MPS backend.

\FloatBarrier
